\documentclass[a4paper]{arrowhead}

\usepackage[yyyymmdd]{datetime}
\usepackage{etoolbox}
\usepackage[utf8]{inputenc}
\usepackage[T1]{fontenc}
\usepackage{multirow}
\usepackage{longtable}
\usepackage{tikz}
\usepackage{adjustbox}
\AtBeginEnvironment{longtable}{\small}
\BeforeBeginEnvironment{tikzpicture}{\begin{adjustbox}{max width=\textwidth}}
\AfterEndEnvironment{tikzpicture}{\end{adjustbox}}
\usetikzlibrary{positioning,shapes.multipart,arrows.meta}

\renewcommand{\dateseparator}{-}
\setlength{\emergencystretch}{3em}
\lstset{language=HTTP,breakatwhitespace=false,columns=fullflexible}

\newcommand{\fparam}[1]{\textit{\textcolor{ArrowheadBlue}{#1}}}

%% Special references
\newcommand{\mref}[2][]{{\textcolor{ArrowheadPurple}{\hyperref[sec:model:#2]{\ifstrempty{#1}{#2}{#1}}}}}
\newcommand{\oref}[1]{{\textcolor{ArrowheadBlue}{\hyperref[sec:operations:#1]{#1}}}}
\newcommand{\pdef}[1]{{\textcolor{ArrowheadGrey}{#1 \label{sec:model:primitives:#1} \label{sec:model:primitives:#1s}}}}
\newcommand{\pref}[1]{{\textcolor{ArrowheadGrey}{\hyperref[sec:model:primitives:#1]{#1}}}}

% as template_idd.tex; the input and output arguments are passed already formatted (\mref or \pref)
\newrobustcmd\osubsection[5]{
  \addtocounter{subsection}{1}
  \addcontentsline{toc}{subsection}{\protect\numberline{\thesubsection}operation \textcolor{ArrowheadBlue}{#1}}
  \subsection*{
    \thesubsection\quad
    #2 \textcolor{ArrowheadPurple}{#3} \\
    \small
    \hspace*{0.075\textwidth}\begin{minipage}{0.1\textwidth}
      \vspace*{1mm}
      Operation: \\
      \notblank{#4}{Input: \\}{}
      \notblank{#5}{Output: \\}{}
    \end{minipage}%
    \begin{minipage}{0.825\textwidth}
      \vspace*{1mm}
      \textcolor{ArrowheadBlue}{#1} \\
      \notblank{#4}{#4 \\}{}
      \notblank{#5}{#5 \\}{}
    \end{minipage}
  }
  \label{sec:operations:#1}
}
\newrobustcmd\msubsection[2]{
  \addtocounter{subsection}{1}
  \addcontentsline{toc}{subsection}{\protect\numberline{\thesubsection}#1 \textcolor{ArrowheadPurple}{#2}}
  \subsection*{\thesubsection\quad#1 \textcolor{ArrowheadPurple}{#2}}
  \label{sec:model:#2}
}
%%

\usepackage{xurl}
\begin{document}
\sloppy

%% Arrowhead Document Properties
\ArrowheadTitle{authorizationTokenManagement MQTT/TLS/JSON}
\ArrowheadServiceID{authorizationTokenManagement}
\ArrowheadType{Interface Design Description}
\ArrowheadTypeShort{IDD}
\ArrowheadVersion{5.3.0}
\ArrowheadDate{2026-11-09}
\ArrowheadAuthor{Jerker Delsing}
\ArrowheadStatus{REVIEW DRAFT}
\ArrowheadContact{jerker.delsing@ltu.se}
\ArrowheadFooter{\href{www.arrowhead.eu}{www.arrowhead.eu}}
\ArrowheadSetup
%%

%% Front Page
\begin{center}
  \vspace*{1cm}
  \huge{\arrowtitle}

  \vspace*{0.2cm}
  \LARGE{\arrowtype}
  \vspace*{1cm}

  \Large{Service ID: \textit{"\arrowid"}}

  \vspace*{\fill}

  % Front Page Image

  \vspace*{1cm}
  \vspace*{\fill}

  % Front Page Abstract
  \begin{abstract}
    This document describes the MQTT protocol with TLS payload security and JSON payload encoding variant of the \texttt{authorization\allowbreak{}Token\allowbreak{}Management} service of the ConsumerAuthorization system.
  \end{abstract}

  \vspace*{1cm}
\end{center}

\newpage
%%

%% Table of Contents
\tableofcontents
\newpage
%%

\section{Documentation edition and implementation scope}
This is the 5.3.0 documentation edition, dated 2026-11-09. It carries forward the system and service contracts from the retained 5.2.2 documentation and identifies the current primary implementation sources. The former \path|documentation/v5.2.2| folder is retained unchanged.

Document version 5.3.0 is distinct from executable, dependency, registered system/service and wire-contract versions. Values still defined by the source (including 5.2.2 artefact versions, 1.0.0 service versions and \path|SysMonitor-AMQP-IDD-5.2.2| payload identifiers) are preserved. This documentation update does not change runtime behaviour or establish a new implementation release. Historical verification statements describe the tests recorded in the inherited baseline; this revision validates document compilation and references.

\section{Overview}
\label{sec:overview}

This document describes the \texttt{authorization\allowbreak{}Token\allowbreak{}Management} service interface of the ConsumerAuthorization system over MQTT, which provides operator management of authorization tokens and of provider encryption keys.

It's implemented using protocol, encoding as stated in below table:
\begin{table}[ht!]
  \centering
  \begin{tabular}{|l|l|l|}
    \hline
    \rowcolor{gray!33} Profile type & Type & Version \\ \hline
    Transfer protocol & MQTT & 3.1.1 \cite{mqtt311} \\ \hline
    Data encryption & TLS & 1.3 \cite{rfc8446}, 1.2 \cite{rfc5246} \\ \hline
    Encoding & JSON & RFC 8259 \cite{rfc8259} \\ \hline
    Compression & N/A & - \\ \hline
    Semantics & N/A & - \\ \hline
    Ontology & N/A & - \\ \hline
  \end{tabular}
  \caption{Communication and semantics details used for the \texttt{authorization\allowbreak{}Token\allowbreak{}Management} service interface}
  \label{tab:comunication_semantics_profile}
\end{table}

This document provides the Interface Design Description IDD to the \textit{authorizationTokenManagement -- Service Description} document \cite{ca-authorizationTokenManagement-sd}.
The operations, their semantics and the data models are the same as in the HTTP/TLS/JSON IDD \cite{ca-authorizationTokenManagement-http-idd}; this document specifies the MQTT binding.

The MQTT interface is offered only when the system runs with \texttt{mqtt.api.enabled=true}; it is then registered with the interface template \texttt{generic\_\allowbreak{}mqtts} (\texttt{generic\_\allowbreak{}mqtt} without TLS), the broker address and port (\texttt{mqtt.\allowbreak{}broker.\allowbreak{}address}, \texttt{mqtt.\allowbreak{}broker.\allowbreak{}port}) and the base topic \texttt{arrowhead/\allowbreak{}consumer-\allowbreak{}authorization/\allowbreak{}authorization-\allowbreak{}token/\allowbreak{}management/\allowbreak{}}. In the Arrowhead-test-cloud-java deployment the MQTT interfaces of the core systems are disabled.

\newpage

\section{Service Operations}
\label{sec:operations}

The provider subscribes to one topic per operation, the base topic followed by the operation name:

\noindent\begin{longtable}{| p{5cm} | p{10.4cm} |} \hline
\rowcolor{gray!33} Operation & Request topic \\ \hline \endhead
\texttt{generate-\allowbreak{}tokens} & \texttt{arrowhead/\allowbreak{}consumer-\allowbreak{}authorization/\allowbreak{}authorization-\allowbreak{}token/\allowbreak{}management/\allowbreak{}generate-\allowbreak{}tokens} \\ \hline
\texttt{query-\allowbreak{}tokens} & \texttt{arrowhead/\allowbreak{}consumer-\allowbreak{}authorization/\allowbreak{}authorization-\allowbreak{}token/\allowbreak{}management/\allowbreak{}query-\allowbreak{}tokens} \\ \hline
\texttt{revoke-\allowbreak{}tokens} & \texttt{arrowhead/\allowbreak{}consumer-\allowbreak{}authorization/\allowbreak{}authorization-\allowbreak{}token/\allowbreak{}management/\allowbreak{}revoke-\allowbreak{}tokens} \\ \hline
\texttt{add-\allowbreak{}encryption-\allowbreak{}keys} & \texttt{arrowhead/\allowbreak{}consumer-\allowbreak{}authorization/\allowbreak{}authorization-\allowbreak{}token/\allowbreak{}management/\allowbreak{}add-\allowbreak{}encryption-\allowbreak{}keys} \\ \hline
\texttt{remove-\allowbreak{}encryption-\allowbreak{}keys} & \texttt{arrowhead/\allowbreak{}consumer-\allowbreak{}authorization/\allowbreak{}authorization-\allowbreak{}token/\allowbreak{}management/\allowbreak{}remove-\allowbreak{}encryption-\allowbreak{}keys} \\ \hline
\end{longtable}

A consumer publishes a request envelope (\mref{MqttRequestTemplate}) to the topic and receives the answer (\mref{MqttResponseTemplate}) on its own \texttt{responseTopic}. The \texttt{authentication} field carries the requester's identity according to the provider's authentication policy: \texttt{IDENTITY-TOKEN//<token>} (outsourced or internal policy), \texttt{SYSTEM//<name>} (declared policy) or the requester's X.509 certificate (certificate policy). The \texttt{status} of the response uses the HTTP status code numbers (200, 201, 204, 400, 401, 403, 404, 408, 423, 500, 503); on error the payload is an \mref{ErrorMessage}.

\begin{lstlisting}[caption={Request envelope.}]
{
  "traceId": "a1b2c3",
  "authentication": "IDENTITY-TOKEN//<token>",
  "responseTopic": "arrowhead/response/TemperatureDisplay/a1b2c3",
  "qosRequirement": 1,
  "params": {},
  "payload": {}
}
\end{lstlisting}

\begin{lstlisting}[caption={Response envelope.}]
{
  "status": 200,
  "traceId": "a1b2c3",
  "receiver": "TemperatureDisplay",
  "payload": {}
}
\end{lstlisting}

\osubsection{generate-tokens}{PUBLISH}{\texttt{arrowhead/\allowbreak{}consumer-\allowbreak{}authorization/\allowbreak{}authorization-\allowbreak{}token/\allowbreak{}management/\allowbreak{}generate-\allowbreak{}tokens}}{\mref[Authorization\allowbreak{}Token\allowbreak{}Generation\allowbreak{}Mgmt\allowbreak{}List\allowbreak{}Request]{AuthorizationTokenGenerationMgmtListRequest}}{\mref[Authorization\allowbreak{}Token\allowbreak{}Mgmt\allowbreak{}List\allowbreak{}Response]{AuthorizationTokenMgmtListResponse}}

Generates tokens on behalf of consumers. With \texttt{unbound=\allowbreak{}true} the authorization check is skipped (only for systems in \texttt{unbounded.\allowbreak{}token.\allowbreak{}generation.\allowbreak{}whitelist}).

Optional entries of \texttt{params}: \texttt{unbound}.

\begin{lstlisting}[caption={Payload of an \oref{generate-tokens} request.}]
{
  "list": [
    {
      "tokenVariant": "TIME_LIMITED_TOKEN_AUTH",
      "targetType": "SERVICE_DEF",
      "consumerCloud": "LOCAL",
      "consumer": "TemperatureDisplay",
      "provider": "TemperatureSensor",
      "target": "temperatureReading",
      "scope": "read",
      "expiresAt": "2026-10-22T12:00:00Z",
      "usageLimit": 5
    }
  ]
}
\end{lstlisting}

\osubsection{query-tokens}{PUBLISH}{\texttt{arrowhead/\allowbreak{}consumer-\allowbreak{}authorization/\allowbreak{}authorization-\allowbreak{}token/\allowbreak{}management/\allowbreak{}query-\allowbreak{}tokens}}{\mref[Authorization\allowbreak{}Token\allowbreak{}Query\allowbreak{}Request]{AuthorizationTokenQueryRequest}}{\mref[Authorization\allowbreak{}Token\allowbreak{}Mgmt\allowbreak{}List\allowbreak{}Response]{AuthorizationTokenMgmtListResponse}}

Returns a page of tokens matching the filter.

\begin{lstlisting}[caption={Payload of an \oref{query-tokens} request.}]
{
  "pagination": {
    "page": 0,
    "size": 10,
    "direction": "ASC",
    "sortField": "createdAt"
  },
  "requester": "DynamicServiceOrchestration",
  "tokenType": "TIME_LIMITED_TOKEN",
  "consumerCloud": "LOCAL",
  "consumer": "TemperatureDisplay",
  "provider": "TemperatureSensor",
  "targetType": "SERVICE_DEF",
  "target": "temperatureReading"
}
\end{lstlisting}

\osubsection{revoke-tokens}{PUBLISH}{\texttt{arrowhead/\allowbreak{}consumer-\allowbreak{}authorization/\allowbreak{}authorization-\allowbreak{}token/\allowbreak{}management/\allowbreak{}revoke-\allowbreak{}tokens}}{\pref{Array}$<$\allowbreak{}\pref{String}$>$}{}

Revokes tokens by token reference.

\begin{lstlisting}[caption={Payload of an \oref{revoke-tokens} request.}]
[
  "<token reference>"
]
\end{lstlisting}

\osubsection{add-encryption-keys}{PUBLISH}{\texttt{arrowhead/\allowbreak{}consumer-\allowbreak{}authorization/\allowbreak{}authorization-\allowbreak{}token/\allowbreak{}management/\allowbreak{}add-\allowbreak{}encryption-\allowbreak{}keys}}{\mref[Authorization\allowbreak{}Mgmt\allowbreak{}Encryption\allowbreak{}Key\allowbreak{}Registration\allowbreak{}List\allowbreak{}Request]{AuthorizationMgmtEncryptionKeyRegistrationListRequest}}{\mref[Authorization\allowbreak{}Mgmt\allowbreak{}Encryption\allowbreak{}Key\allowbreak{}List\allowbreak{}Response]{AuthorizationMgmtEncryptionKeyListResponse}}

Registers encryption keys on behalf of providers.

\begin{lstlisting}[caption={Payload of an \oref{add-encryption-keys} request.}]
{
  "list": [
    {
      "systemName": "TemperatureSensor",
      "key": "0123456789abcdef",
      "algorithm": "AES/ECB/PKCS5Padding"
    }
  ]
}
\end{lstlisting}

\osubsection{remove-encryption-keys}{PUBLISH}{\texttt{arrowhead/\allowbreak{}consumer-\allowbreak{}authorization/\allowbreak{}authorization-\allowbreak{}token/\allowbreak{}management/\allowbreak{}remove-\allowbreak{}encryption-\allowbreak{}keys}}{\pref{Array}$<$\allowbreak{}\pref{String}$>$}{}

Removes the encryption keys of the named providers.

\begin{lstlisting}[caption={Payload of an \oref{remove-encryption-keys} request.}]
[
  "TemperatureSensor"
]
\end{lstlisting}

\newpage

\section{Data Models}
\label{sec:model}

Here, all data objects that can be part of the service calls associated with this service are listed in alphabetic order.
Each \textit{struct} denotes a JSON \pref{Object} with the named fields; fields that are not needed may be omitted.

\msubsection{struct}{MqttRequestTemplate}

The request envelope.

\noindent\begin{tabularx}{\textwidth}{| p{4.25cm} | p{3.5cm} | X |} \hline
\rowcolor{gray!33} Object Field & Value Type & Description \\ \hline
"traceId" & \pref{String} & Identifier chosen by the requester, returned in the response. \\ \hline
"authentication" & \pref{String} & Identity of the requester (see Section \ref{sec:operations}). \\ \hline
"responseTopic" & \pref{String} & Topic on which the response is published. \\ \hline
"qosRequirement" & \pref{Number} & MQTT QoS of the response (0, 1 or 2). \\ \hline
"params" & \pref{Object}$<$\pref{String}$>$ & Optional parameters (the query parameters of the HTTP binding). \\ \hline
"payload" & \pref{Object} & The operation input. \\ \hline
\end{tabularx}

\msubsection{struct}{MqttResponseTemplate}

The response envelope.

\noindent\begin{tabularx}{\textwidth}{| p{4.25cm} | p{3.5cm} | X |} \hline
\rowcolor{gray!33} Object Field & Value Type & Description \\ \hline
"status" & \pref{Number} & Status code. \\ \hline
"traceId" & \pref{String} & The trace identifier of the request. \\ \hline
"receiver" & \pref{String} & Name of the requester. \\ \hline
"payload" & \pref{Object} & The operation output, or an \mref{ErrorMessage}. \\ \hline
\end{tabularx}

\msubsection{struct}{ErrorMessage}




\noindent\begin{tabularx}{\textwidth}{| >{\raggedright\arraybackslash}p{3.9cm} | >{\raggedright\arraybackslash\small}p{5.3cm} | X |} \hline
\rowcolor{gray!33} Object Field & Value Type & Description \\ \hline
"error\allowbreak{}Message" & \pref{String} &  \\ \hline
"error\allowbreak{}Code" & \pref{Number} &  \\ \hline
"exception\allowbreak{}Type" & \pref{ExceptionType} &  \\ \hline
"origin" & \pref{String} &  \\ \hline
\end{tabularx}



\msubsection{struct}{AuthorizationMgmtEncryptionKeyListResponse}

List of registered encryption keys.


\noindent\begin{tabularx}{\textwidth}{| >{\raggedright\arraybackslash}p{3.9cm} | >{\raggedright\arraybackslash\small}p{5.3cm} | X |} \hline
\rowcolor{gray!33} Object Field & Value Type & Description \\ \hline
"entries" & \pref{Array}$<$\allowbreak{}\mref[Authorization\allowbreak{}Mgmt\allowbreak{}Encryption\allowbreak{}Key\allowbreak{}Response]{AuthorizationMgmtEncryptionKeyResponse}$>$ & Matching records of the current page. \\ \hline
"count" & \pref{Number} & Total number of matching records (all pages). \\ \hline
\end{tabularx}



\msubsection{struct}{AuthorizationMgmtEncryptionKeyRegistrationListRequest}

Encryption keys to register on behalf of providers.


\noindent\begin{tabularx}{\textwidth}{| >{\raggedright\arraybackslash}p{3.9cm} | >{\raggedright\arraybackslash\small}p{5.3cm} | X |} \hline
\rowcolor{gray!33} Object Field & Value Type & Description \\ \hline
"list" & \pref{Array}$<$\allowbreak{}\mref[Authorization\allowbreak{}Mgmt\allowbreak{}Encryption\allowbreak{}Key\allowbreak{}Registration\allowbreak{}Request]{AuthorizationMgmtEncryptionKeyRegistrationRequest}$>$ & The list of request items. \\ \hline
\end{tabularx}



\msubsection{struct}{AuthorizationMgmtEncryptionKeyRegistrationRequest}

Encryption key of one provider.


\noindent\begin{tabularx}{\textwidth}{| >{\raggedright\arraybackslash}p{3.9cm} | >{\raggedright\arraybackslash\small}p{5.3cm} | X |} \hline
\rowcolor{gray!33} Object Field & Value Type & Description \\ \hline
"system\allowbreak{}Name" & \pref{String} & Provider system the key belongs to. \\ \hline
"key" & \pref{String} & The encryption key, as a string. \\ \hline
"algorithm" & \pref{String} & Encryption algorithm of the key, e.g. \texttt{AES/\allowbreak{}ECB/\allowbreak{}PKCS5Padding} or \texttt{AES/\allowbreak{}CBC/\allowbreak{}PKCS5Padding}. \\ \hline
\end{tabularx}



\msubsection{struct}{AuthorizationMgmtEncryptionKeyResponse}

A registered encryption key.


\noindent\begin{tabularx}{\textwidth}{| >{\raggedright\arraybackslash}p{3.9cm} | >{\raggedright\arraybackslash\small}p{5.3cm} | X |} \hline
\rowcolor{gray!33} Object Field & Value Type & Description \\ \hline
"system\allowbreak{}Name" & \pref{String} & Name of the system, PascalCase (e.g. \texttt{TemperatureSensor}). \\ \hline
"raw\allowbreak{}Key" & \pref{String} & The registered encryption key. \\ \hline
"algorithm" & \pref{String} & Encryption algorithm of the key, e.g. \texttt{AES/\allowbreak{}ECB/\allowbreak{}PKCS5Padding} or \texttt{AES/\allowbreak{}CBC/\allowbreak{}PKCS5Padding}. \\ \hline
"key\allowbreak{}Additive" & \pref{String} & Initialization vector (CBC mode), Base64; empty for ECB. \\ \hline
"created\allowbreak{}At" & \pref{DateTime} & UTC time of creation. \\ \hline
\end{tabularx}



\msubsection{struct}{AuthorizationTokenGenerationMgmtListRequest}

Tokens to generate on behalf of consumers.


\noindent\begin{tabularx}{\textwidth}{| >{\raggedright\arraybackslash}p{3.9cm} | >{\raggedright\arraybackslash\small}p{5.3cm} | X |} \hline
\rowcolor{gray!33} Object Field & Value Type & Description \\ \hline
"list" & \pref{Array}$<$\allowbreak{}\mref[Authorization\allowbreak{}Token\allowbreak{}Generation\allowbreak{}Mgmt\allowbreak{}Request]{AuthorizationTokenGenerationMgmtRequest}$>$ & The list of request items. \\ \hline
\end{tabularx}



\msubsection{struct}{AuthorizationTokenGenerationMgmtRequest}

One token to generate on behalf of a consumer.


\noindent\begin{tabularx}{\textwidth}{| >{\raggedright\arraybackslash}p{3.9cm} | >{\raggedright\arraybackslash\small}p{5.3cm} | X |} \hline
\rowcolor{gray!33} Object Field & Value Type & Description \\ \hline
"token\allowbreak{}Variant" & \pref{String} & Requested token variant, as interface security policy (e.g. \texttt{TIME\_\allowbreak{}LIMITED\_\allowbreak{}TOKEN\_\allowbreak{}AUTH}, \texttt{RSA\_\allowbreak{}SHA512\_\allowbreak{}JSON\_\allowbreak{}WEB\_\allowbreak{}TOKEN\_\allowbreak{}AUTH}). \\ \hline
"target\allowbreak{}Type" & \pref{String} & Target type: \texttt{SERVICE\_\allowbreak{}DEF} or \texttt{EVENT\_\allowbreak{}TYPE}. \\ \hline
"consumer\allowbreak{}Cloud" & \pref{String} & Cloud identifier of the consumer system. \\ \hline
"consumer" & \pref{String} & Name of the consumer system. \\ \hline
"provider" & \pref{String} & Name of the provider system. \\ \hline
"target" & \pref{String} & Target of the policy or token: a service definition or event type. \\ \hline
"scope" & \pref{String} & Operation the token or policy is limited to; empty for all operations. \\ \hline
"expires\allowbreak{}At" & \pref{DateTime} & Requested expiry; default \texttt{token.\allowbreak{}time.\allowbreak{}limit} from now. \\ \hline
"usage\allowbreak{}Limit" & \pref{Number} & Requested uses; default \texttt{simple.\allowbreak{}token.\allowbreak{}usage.\allowbreak{}limit}. \\ \hline
\end{tabularx}



\msubsection{struct}{AuthorizationTokenMgmtListResponse}

List of generated tokens.


\noindent\begin{tabularx}{\textwidth}{| >{\raggedright\arraybackslash}p{3.9cm} | >{\raggedright\arraybackslash\small}p{5.3cm} | X |} \hline
\rowcolor{gray!33} Object Field & Value Type & Description \\ \hline
"entries" & \pref{Array}$<$\allowbreak{}\mref[Authorization\allowbreak{}Token\allowbreak{}Response]{AuthorizationTokenResponse}$>$ & Matching records of the current page. \\ \hline
"count" & \pref{Number} & Total number of matching records (all pages). \\ \hline
\end{tabularx}



\msubsection{struct}{AuthorizationTokenQueryRequest}

Filter for generated tokens.


\noindent\begin{tabularx}{\textwidth}{| >{\raggedright\arraybackslash}p{3.9cm} | >{\raggedright\arraybackslash\small}p{5.3cm} | X |} \hline
\rowcolor{gray!33} Object Field & Value Type & Description \\ \hline
"pagination" & \mref{Page} & Paging and sorting of the result. \\ \hline
"requester" & \pref{String} & Name of the system that requested the token. \\ \hline
"token\allowbreak{}Type" & \pref{String} & Token type to match. \\ \hline
"consumer\allowbreak{}Cloud" & \pref{String} & Cloud identifier of the consumer system. \\ \hline
"consumer" & \pref{String} & Name of the consumer system. \\ \hline
"provider" & \pref{String} & Name of the provider system. \\ \hline
"target\allowbreak{}Type" & \pref{String} & Target type: \texttt{SERVICE\_\allowbreak{}DEF} or \texttt{EVENT\_\allowbreak{}TYPE}. \\ \hline
"target" & \pref{String} & Target of the policy or token: a service definition or event type. \\ \hline
\end{tabularx}



\msubsection{struct}{AuthorizationTokenResponse}

A generated token with its bookkeeping data.


\noindent\begin{tabularx}{\textwidth}{| >{\raggedright\arraybackslash}p{3.9cm} | >{\raggedright\arraybackslash\small}p{5.3cm} | X |} \hline
\rowcolor{gray!33} Object Field & Value Type & Description \\ \hline
"token\allowbreak{}Type" & \pref{AuthorizationTokenType} & Token type: \texttt{TIME\_\allowbreak{}LIMITED\_\allowbreak{}TOKEN}, \texttt{USAGE\_\allowbreak{}LIMITED\_\allowbreak{}TOKEN}, \texttt{SELF\_\allowbreak{}CONTAINED\_\allowbreak{}TOKEN} or \texttt{TRANSLATION\_\allowbreak{}BRIDGE\_\allowbreak{}TOKEN}. \\ \hline
"variant" & \pref{String} & Token variant (interface security policy). \\ \hline
"token" & \pref{String} & The token, for simple tokens only in hashed form. \\ \hline
"token\allowbreak{}Reference" & \pref{String} & Reference of the token, used to revoke it. \\ \hline
"requester" & \pref{String} & Name of the system that requested the token. \\ \hline
"consumer\allowbreak{}Cloud" & \pref{String} & Cloud identifier of the consumer system. \\ \hline
"consumer" & \pref{String} & Name of the consumer system. \\ \hline
"provider" & \pref{String} & Name of the provider system. \\ \hline
"target\allowbreak{}Type" & \pref{AuthorizationTargetType} & Target type: \texttt{SERVICE\_\allowbreak{}DEF} or \texttt{EVENT\_\allowbreak{}TYPE}. \\ \hline
"target" & \pref{String} & Target of the policy or token: a service definition or event type. \\ \hline
"scope" & \pref{String} & Operation the token or policy is limited to; empty for all operations. \\ \hline
"created\allowbreak{}At" & \pref{DateTime} & UTC time of creation. \\ \hline
"usage\allowbreak{}Limit" & \pref{Number} & Number of allowed uses (usage-limited tokens). \\ \hline
"usage\allowbreak{}Left" & \pref{Number} & Remaining uses of a usage-limited token. \\ \hline
"expires\allowbreak{}At" & \pref{DateTime} & UTC expiry time; no value means no expiry. \\ \hline
\end{tabularx}



\msubsection{struct}{Page}

Paging and sorting; all fields optional.


\noindent\begin{tabularx}{\textwidth}{| >{\raggedright\arraybackslash}p{3.9cm} | >{\raggedright\arraybackslash\small}p{5.3cm} | X |} \hline
\rowcolor{gray!33} Object Field & Value Type & Description \\ \hline
"page" & \pref{Number} & Zero-based page number. \\ \hline
"size" & \pref{Number} & Number of records per page (at most \texttt{max.\allowbreak{}page.\allowbreak{}size}, default 1000). \\ \hline
"direction" & \pref{String} & Sort direction, \texttt{ASC} or \texttt{DESC}. \\ \hline
"sort\allowbreak{}Field" & \pref{String} & Field to sort by. \\ \hline
\end{tabularx}



\subsection{Primitives}
\label{sec:model:primitives}

As all messages are encoded using the JSON format \cite{rfc8259}, the following primitive constructs, part of that standard, become available.

\noindent\begin{tabularx}{\textwidth}{| p{4.6cm} | X |} \hline
\rowcolor{gray!33} Type & Description \\ \hline
\pdef{Array} & A JSON Array \cite{rfc8259}; \pref{Array}$<$A$>$ has elements of type A. \\ \hline
\pdef{AuthorizationTargetType} & One of \texttt{SERVICE\_\allowbreak{}DEF}, \texttt{EVENT\_\allowbreak{}TYPE} (JSON String). \\ \hline
\pdef{AuthorizationTokenType} & One of \texttt{TIME\_\allowbreak{}LIMITED\_\allowbreak{}TOKEN}, \texttt{USAGE\_\allowbreak{}LIMITED\_\allowbreak{}TOKEN}, \texttt{SELF\_\allowbreak{}CONTAINED\_\allowbreak{}TOKEN}, \texttt{TRANSLATION\_\allowbreak{}BRIDGE\_\allowbreak{}TOKEN} (JSON String). \\ \hline
\pdef{DateTime} & JSON String with an ISO 8601 UTC timestamp, e.g. \texttt{2026-10-22T12:00:00Z}. \\ \hline
\pdef{ExceptionType} & One of \texttt{ARROWHEAD}, \texttt{INVALID\_\allowbreak{}PARAMETER}, \texttt{AUTH}, \texttt{FORBIDDEN}, \texttt{DATA\_\allowbreak{}NOT\_\allowbreak{}FOUND}, \texttt{TIMEOUT}, \texttt{LOCKED}, \texttt{INTERNAL\_\allowbreak{}SERVER\_\allowbreak{}ERROR}, \texttt{EXTERNAL\_\allowbreak{}SERVER\_\allowbreak{}ERROR} (JSON String). \\ \hline
\pdef{Number} & JSON Number, an integer. \\ \hline
\pdef{Object} & JSON Object; \pref{Object}$<$A$>$ has string keys and values of type A. \\ \hline
\pdef{String} & JSON String. \\ \hline
\end{tabularx}

\newpage

\section{Go and Rust implementation variants}
\subsection{Go implementation}
Primary system folder: \path|Eclipse-Arrowhead-go/ah5-core-go/consumer-authorization|.\par
The system owns its Go module, launcher and build configuration. Use the per-system \path|Makefile| and \path|go.mod|; the containing Go tree is the build source, not a historical Arrowhead copy. Ordinary native roles use independent identities and JSON snapshot state rather than the Java database. Do not assume Java SQL migration, clustering or identical transaction/retry semantics. The registry has its own configuration and optional native AMQP monitor; other roles must be checked against their own registrations.
\paragraph{Listening port.} Built-in default 127.0.0.1:8445 (Arrowhead default). The shipped config.json sets "listen": "0.0.0.0:8445". Override with "listen" in the configuration file, SERVER\_PORT=N or -server.port N; the port overrides keep the configured host.
Native design/route evidence: \path|Eclipse-Arrowhead-go/ah5-core-go/consumer-authorization/documentation|.\par
For this Java interface, select only a native transport and operation explicitly implemented and advertised by that port. A matching method/path does not establish DTO, error, authorization, retry or persistence equivalence. Java MQTT/AMQP bindings are not automatically provided by a Go or Rust role. Differential contract and security checks remain necessary before substitution.
\subsection{Rust implementation}
Primary system folder: \path|Eclipse-Arrowhead-rust/ah5-core-rust/consumer-authorization|.\par
The per-system entry point is built by the \path|arrowhead-core-rust| crate. Build from the primary Rust workspace with the applicable package/binary selection. Documentation links can resolve into the shared core/support catalogue. Independent process identity and state remain required.\par
The ordinary native profile uses Tokio/Axum and rustls with peer-identity checks and independent locked JSON state. Network effects are not distributed database transactions. The precise trust, CRL reload and persistence guarantees are those of the native profile, not the Java implementation.\par
Native design/route evidence: \path|Eclipse-Arrowhead-rust/ah5-core-rust/documentation/ConsumerAuthorization|.\par
For this Java interface, select only a native transport and operation explicitly implemented and advertised by that port. A matching method/path does not establish DTO, error, authorization, retry or persistence equivalence. Java MQTT/AMQP bindings are not automatically provided by a Go or Rust role. Differential contract and security checks remain necessary before substitution.

\bibliographystyle{IEEEtran}
\bibliography{bibliography}

\newpage

\section{Revision History}
\subsection{Amendments}

\noindent\begin{tabularx}{\textwidth}{| p{1cm} | p{3cm} | p{2cm} | X | p{4cm} |} \hline
\rowcolor{gray!33} No. & Date & Version & Subject of Amendments & Author \\ \hline

1 & 2026-10-22 & 5.2.2 & Document created from the Arrowhead LaTeX templates for the 5.2.2 Java Spring implementation, including the local extensions of this workspace. & Jerker Delsing \\ \hline
2 & 2026-11-09 & 5.3.0 & New documentation edition; primary-source references and Go/Rust implementation scope; 5.2.2 retained. & Documentation maintenance \\ \hline
\end{tabularx}

\subsection{Quality Assurance}

\noindent\begin{tabularx}{\textwidth}{| p{1cm} | p{3cm} | p{2cm} | X |} \hline
\rowcolor{gray!33} No. & Date & Version & Approved by \\ \hline

1 & & \arrowversion  &  \\ \hline

\end{tabularx}

\end{document}
